\documentclass[11pt,a4paper]{article}
\usepackage[T1]{fontenc}
\usepackage[utf8]{inputenc}
\usepackage[spanish,es-nodecimaldot]{babel}
\usepackage{amsmath,amssymb}
\usepackage{booktabs}
\usepackage{graphicx}
\usepackage{siunitx}
\usepackage[margin=2.5cm]{geometry}
\usepackage[hidelinks]{hyperref}

\title{Laboratorio 1: Python Científico Legible y Reproducible}
\author{Integrantes: Cristian Stiven Capera, Ana Sofia García, Jeimmy Andrea Gonzáles, Karol Mariana Gutiérrez y Brayan Andrés Ruiz}
\date{\today}

\begin{document}
\maketitle

\begin{abstract}
Este informe muestra el desarrollo de un proyecto en Python que puede ejecutarse en cualquier computador y obtener exactamente los mismos resultados. La idea principal es comprobar la replicación de un análisis completo (con datos sintéticos, estadísticas, validaciones y gráficos) sin necesidad de copiar entornos virtuales, importar nuevas librerías, ni cambiar rutas a mano. Para lograrlo, se organiza el código por módulos y se usan herramientas como \texttt{uv} para las dependencias, \texttt{ruff} para mantener el código limpio, \texttt{pytest} para las pruebas automáticas y LaTeX para generar documentos de forma automática.
\end{abstract}

\section{Identificación y Pregunta de Trabajo}
La consolidación del laboratorio fue guiada por una pregunta base, que es: \textbf{¿puede el equipo reconstruir en un computador distinto un análisis que conserve los mismos parámetros, datos sintéticos, estadísticas, validaciones y figuras, sin copiar el entorno virtual ni editar rutas?}

\section{Arquitectura del Proyecto}
El proyecto se organiza bajo una separación de responsabilidades, garantizando que los componentes se relacionen de manera modular. Para entender la distribución física del repositorio, la estructura general del proyecto se divide de la siguiente manera:

\begin{verbatim}
equipo_02_lab01_python_reproducible/
|-- .python-version
|-- pyproject.toml
|-- uv.lock
|-- README.md
|-- data/
|   |-- raw/
|   `-- processed/
|-- src/
|   `-- python_scientific/
|       |-- __init__.py
|       |-- quality.py
|       |-- descriptive_statistics.py
|       |-- didactic_transformation.py
|       |-- signal_derivation.py
|       |-- signal_integration.py
|       `-- three_panel_figure.py
|-- tests/
|   `-- test_pipeline.py
|-- scripts/
|   `-- reproduce.py
|-- results/
|   |-- algebra_report.json
|   |-- contributions.md
|   |-- delivery_manifest.txt
|   |-- derivation_integration_report.json
|   |-- quality_report.json
|   |-- stats.json
|   `-- summary_table.tex
|-- figures/
|   `-- summary.png
`-- reports/
    `-- informe.tex
\end{verbatim}

\noindent\ A partir de esta organización, la Tabla~\ref{tab:auditoria} detalla los flujos de entrada, transformación, salida y la capacidad de regeneración de cada elemento clave del repositorio.

\begin{table}[htbp]
  \centering
  \caption{Auditoría de estructura, responsabilidades y regeneración del proyecto.}
  \label{tab:auditoria}
  \small
  \begin{tabular}{p{2.5cm} p{3.2cm} p{3.2cm} p{3.2cm} c}
    \toprule
    \textbf{Archivo / Carpeta} & \textbf{Entrada} & \textbf{Transformación} & \textbf{Salida} & \textbf{Regenerable} \\
    \midrule
    \texttt{pyproject.toml} & Definición de metadatos y dependencias & Configura el entorno y las reglas de formato (\texttt{ruff}) & Configuración estandarizada del proyecto & Sí \\
    \addlinespace
    \texttt{uv.lock} & Versiones de paquetes declaradas & Bloquea las versiones milimétricas exactas de las dependencias & Entorno determinista para sincronización & Sí \\
    \addlinespace
    \texttt{src/} & Funciones de lógica y parámetros base & Aplica modelos numéricos, validaciones de canales y procesamiento & Módulos reutilizables para scripts y pruebas & No \\
    \addlinespace
    \texttt{tests/} & Código fuente en \texttt{src/} y datos sintéticos de prueba & Ejecuta validaciones de contratos vía \texttt{pytest} & Informes de estado de pruebas unitarias & Sí \\
    \addlinespace
    \texttt{data/raw/} & Entradas externas o parámetros iniciales & Almacena las mediciones crudas base para el pipeline & Datos de entrada listos para consumo del sistema & No \\
    \addlinespace
    \texttt{data/processed/} & Datos crudos y parámetros de limpieza & Limpia, normaliza y ajusta los datos en bruto & Series procesadas para análisis estadístico y matemático & Sí \\
    \addlinespace
    \texttt{results/} & Datos procesados y scripts de cálculo & Computa estadísticas, matrices, derivadas e integrales & Tablas formateadas en \LaTeX\ (\texttt{.summary\_table.tex}) y JSON & Sí \\
    \addlinespace
    \texttt{figures/} & Arreglos numéricos de resultados & Genera gráficos analíticos multipanel mediante \texttt{matplotlib} & Archivos de imagen de alta resolución (\texttt{.png}) & Sí \\
    \addlinespace
    \texttt{reports/} & Estructura de texto y salidas de \texttt{results/} y \texttt{figures/} & Compila el documento técnico integrando tablas y figuras dinámicas & Informes y presentaciones en formato PDF final & Sí \\
    \bottomrule
  \end{tabular}
\end{table}


\noindent\ Respecto a la relación entre componentes, en resumen, los scripts en \texttt{scripts/} usan el código de \texttt{src/}, procesan los datos de \texttt{data/} y guardan automáticamente los resultados y las gráficas en \texttt{results/} y \texttt{figures/}. Por último, el informe en \texttt{reports/} toma esos archivos generados para armar el PDF final sin que tengamos que mover nada a mano, asegurando que todo funcione siempre igual.


\section{Métodos}
\subsection{Datos sintéticos, semilla y fórmulas}
Los experimentos numéricos emplean señales sintéticas controladas por una semilla determinista ($\text{seed} = 42$) con un paso temporal estricto $\Delta t = \si{0,05 \second}$. La amplitud de las señales se mide en milivoltios ($\si{\milli\volt}$) y el tiempo en segundos ($\si{\second}$). Para el análisis de mezcla de canales, se implementa la transformación matricial:
\begin{equation}
    \mathbf{y}[n] = \mathbf{A} \mathbf{x}[n]
    \label{eq:mezcla}
\end{equation}

\noindent \textbf{Parámetros explícitos de entrada:}
\begin{itemize}
    \item \textbf{Control de calidad:} Se evaluó un conjunto observado de canales de prueba, para probar la consistencia de los elementos. Dado por \verb|observed = ["SYNTH_A", "SYNTH_B", "SYNTH_A"]| frente al conjunto requerido \verb|required = {"SYNTH_A", "SYNTH_B", "REFERENCE"}|.
    \item \textbf{Señal analizada:} La señal inicial se creó con una fórmula matemática sencilla que mezcla una onda limpia y ruido aleatorio. $x_{\text{raw}}(t) = \sin(2\pi \cdot 0.5 t) + 0.2 \cdot \mathcal{N}(0,1)$. Aquí, la primera parte es una onda senoidal pura con una frecuencia de $0.5$ Hz, y la segunda parte usa ese $\mathcal{N}(0,1)$, que representa un generador de ruido aleatorio (con media cero y desviación uno), multiplicado por $0.2$ para simular las imperfecciones que tendría una medición en la vida real. A partir de esta, también se genera una versión centrada a la que se le sustrae la media.
    \item \textbf{Álgebra lineal:} Se utilizó la matriz de mezcla cuadrada y bien condicionada:
    \begin{equation}
        \mathbf{A} = \begin{pmatrix} 2.0 & 1.0 \\ 1.0 & 3.0 \end{pmatrix}
        \label{eq:matriz_A}
    \end{equation}
    junto con dos canales acoplados ($c_1(t) = \sin(2\pi \cdot 0.05 t)$ y $c_2(t) = \cos(2\pi \cdot 0.1 t)$) para verificar la transformación descrita por:
    \begin{equation}
        \mathbf{y}[n] = \mathbf{A} \mathbf{x}[n]
        \label{eq:mezcla}
    \end{equation}
\end{itemize}

\section{Validación}
\subsection{Control de Calidad}
Para garantizar la integridad estructural de los datos de entrada, se implementó en \texttt{src/quality.py} la función \texttt{validate\_channel\_configuration}. Esta rutina procesa listas de canales observados frente a un conjunto requerido, empleando listas para preservar el orden y conjuntos para detectar faltantes o elementos adicionales. El sistema rechaza estrictamente nombres vacíos y valida casos límite sin mutar los datos originales de entrada.

\subsection{Estadística}
Se evalúan las métricas de tendencia central y dispersión tanto para la señal cruda como para la señal centrada (restando su media), calculando el número de muestras, media, mediana, desviación estándar muestral, mínimo, máximo, rango intercuartílico (IQR) y valor eficaz (RMS). Estos valores se almacenan estructuradamente con sus respectivas unidades en archivos de intercambio.

En cuanto al análisis conceptual de los resultados obtenidos:
\begin{itemize}
    \item \textbf{Media de la señal centrada:} La señal centrada debe tener una media teóricamente cercana a cero (del orden de $10^{-15}$ o menor por limitaciones de precisión en punto flotante), pero no necesariamente un cero absoluto estricto debido al error de redondeo numérico inherente a la representación binaria en la computadora.
    \item \textbf{RMS y desviación estándar:} El valor eficaz (RMS) y la desviación estándar muestral no son idénticos cuando la media de la señal es diferente de cero. Matemáticamente, el RMS incluye tanto la dispersión en torno a la media (varianza) como la energía aportada por el valor medio al cuadrado ($\text{RMS}^2 = \mu^2 + \sigma^2$). Solo coinciden exactamente si la media de la señal es cero.
    \item \textbf{Convención de la desviación estándar:} Para el cálculo de la desviación estándar muestral se empleó la convención con grados de libertad corregidos ($ddof = 1$, dividiendo por $N-1$), lo que proporciona un estimador insesgado de la varianza poblacional a partir de la muestra.
\end{itemize}

\subsection{Álgebra lineal y sistemas de canales}
Para evaluar la recuperación de canales, se define una matriz de mezcla didáctica $\mathbf{A}$. Esta matriz fue seleccionada por ser simétrica y presentar un número de condición bajo, lo que garantiza estabilidad numérica y evita problemas de amplificación de errores.

Tras aplicar la transformación a los canales sintéticos, se realizaron las siguientes etapas metodológicas:
\begin{itemize}
    \item \textbf{Verificación de formas:} Se comprobaron las dimensiones de los arreglos matriciales para asegurar la compatibilidad algebraica en la multiplicación y resolución del sistema.
    \item \textbf{Número de condición:} Se calculó el número de condición de la matriz $\mathbf{A}$, confirmando que se encuentra en un rango bien condicionado y descartando matrices singulares.
    \item \textbf{Recuperación y error:} Utilizando la función \texttt{np.linalg.solve}, se recuperaron los canales originales bajo una tolerancia explícita de $1\times 10^{-7}$, obteniendo un error RMS de reconstrucción despreciable.
\end{itemize}

\textbf{Justificación y efecto de un mal condicionamiento:} 
Si se hubiera empleado una matriz mal condicionada (con un número de condición extremadamente alto), cualquier mínima perturbación numérica o error de redondeo en los datos de entrada se habría amplificado drásticamente al invertir el sistema. Esto habría provocado una pérdida total de estabilidad, derivando en un error RMS elevado y haciendo que los canales recuperados perdieran su validez física.

\subsection{Derivada e Integración}
La derivada temporal se contrasta utilizando como función analítica de referencia:
\begin{equation}
    f(t) = \sin(2\pi f_0 t)
    \label{eq:derivada}
\end{equation}
reportando el error máximo interior y el error RMS en $\si{\milli\volt\per\second}$. En cuanto al tratamiento de extremos, se implementan diferencias finitas hacia adelante y hacia atrás en los bordes para evitar la pérdida de puntos.

Para la integración, se valida la regla trapezoidal mediante la integral polinómica de referencia:
\begin{equation}
    \int_{0}^{1} t^2 \, dt = \frac{1}{3}
    \label{eq:integral}
\end{equation}
aplicada posteriormente de forma acumulada sobre la señal centrada en $\si{\milli\volt\second}$. Se considera una condición inicial nula ($t=0$). El uso de la señal centrada es fundamental para evitar una tendencia lineal ascendente o descendente (rampa) provocada por una línea base no nula.

\section{Resultados y Discusión}
En esta sección se presentan los resultados obtenidos tras la ejecución del pipeline automatizado de procesamiento. Todas las métricas científicas han sido calculadas de manera dinámica a partir de los datos sintéticos generados en el entorno de pruebas. 

\subsection{Evaluación de Control de Calidad}
Como paso inicial, el módulo de validación evaluó la integridad de los canales de la señal mediante el análisis de las estructuras de entrada. El informe generado arrojó los siguientes resultados:
\begin{itemize}
    \item \textbf{Canales duplicados:} Se detectó y reportó el canal \texttt{SYNTH\_A}, identificando solapamientos en la fuente de datos.
    \item \textbf{Canales faltantes:} Se registró la ausencia del canal esperado \texttt{REFERENCE}, lo que activó las alertas preventivas del sistema.
    \item \textbf{Canales extra:} No se encontraron canales adicionales fuera de los esperados (\texttt{extra\_channels: []}).
\end{itemize}

\subsection{Análisis Estadístico de las Señales}
El comportamiento de la señal antes y después del preprocesamiento se resume en la Tabla~\ref{tab:resultados}.

\begin{table}[htbp]
  \centering
  \caption{Resumen estadístico generado de forma automática por el pipeline.}
  \label{tab:resultados}
  \begin{tabular}{lcc}
\toprule
\textbf{Métrica} & \textbf{Señal Cruda (mV)} & \textbf{Señal Centrada (mV)} \\
\midrule
Media & 0.1489 & 2.1316e-17 \\
Mediana & 0.1471 & -0.0019 \\
Desv. Estándar ($ddof=1$) & 0.7329 & 0.7329 \\
IQR & 1.3267 & 1.3267 \\
RMS & 0.7471 & 0.7321 \\
Mínimo / Máximo & -1.20 / 1.60 & -1.35 / 1.45 \\
\bottomrule
\end{tabular}

\end{table}

\subsection{Transformación Lineal y Recuperación de Canales}
El sistema lineal evaluado mediante la matriz de mezcla $\mathbf{A}$ de orden 2 reportó un número de condición adecuado (2.61), lo que permitió que la función \texttt{np.linalg.solve} recuperara los canales originales con un error RMS de reconstrucción por debajo de la tolerancia estricta establecida ($1\times 10^{-7}$). Esto garantiza que la transformación no amplificara errores numéricos de forma indeseada.

\subsection{Derivada e Integración Numérica}
La aplicación de la derivada temporal mediante diferencias finitas y la integración basada en la regla trapezoidal permitieron evaluar el comportamiento dinámico de la señal:
\begin{itemize}
    \item \textbf{Sensibilidad de la derivada:} La derivada obtenida muestra una alta sensibilidad a las pequeñas fluctuaciones de alta frecuencia (ruido estocástico) presentes en la señal sintética, lo que amplifica visualmente las pendientes locales en el tercer panel de la figura 1.
    \item \textbf{Validación de integración:} La validación de la regla trapezoidal implementada devolvió un estado de aprobación (\texttt{regla\_trapecio\_valida: true}), confirmando que el cálculo del área bajo la curva sobre la función de prueba analítica cumple con los umbrales de tolerancia esperados. Así mismo, se registró el valor de la integración acumulada final de la señal centrada.
\end{itemize}

\subsection{Visualización Consolidada}
La Figura~\ref{fig:summary} presenta la visualización de tres paneles generada automáticamente por el pipeline, permitiendo auditar visualmente el impacto de la normalización por media:
\begin{itemize}
    \item \textbf{Panel A:} Presenta la superposición de la señal cruda (en gris) y la señal centrada (en azul).
    \item \textbf{Panel B:} Muestra el histograma de distribución de amplitudes de la señal centrada, el cual permite validar la simetría de los datos respecto al cero.
    \item \textbf{Panel C:} Ilustra la derivada numérica de la señal centrada mediante diferencias finitas en el intervalo de tiempo seleccionado ($1.0$s a $2.0$s), resaltando la dinámica de las variaciones locales.
\end{itemize}

\begin{figure}[h!]
    \centering
    \includegraphics[width=0.85\textwidth]{figures/summary.png}
    \caption{Panel resumen que incluye: (A) Comparativa de señal cruda y centrada, (B) Histograma de amplitudes de la señal centrada, y (C) Derivada numérica en el intervalo $[1.0, 2.0]$ segundos.}
    \label{fig:summary}
\end{figure}

\section{Pruebas}
El proyecto cuenta con un proceso automatizado de 10 pruebas unitarias ejecutadas mediante \texttt{pytest}, que validan rigurosamente los contratos de entrada y salida de cada módulo. El esquema de pruebas cubre:
\begin{itemize}
        \item \textbf{Configuraciones y Calidad (2):} Comprobación de listas completas de canales, detección automática de canales faltantes o duplicados, y rechazo estricto de nombres vacíos sin alterar ni mutar las estructuras de datos originales de entrada.
    \item \textbf{Estadística (3):} Validación con arreglos numéricos conocidos, casos con un solo valor, arreglos vacíos, manejo de valores no finitos (\texttt{NaN}, \texttt{Inf}) y control de formas incorrectas.
    \item \textbf{Álgebra Lineal (1):} Comprueba la correcta resolución de sistemas matriciales, asegurando la consistencia en el cálculo de transformaciones didácticas y validando el control de excepciones ante matrices o vectores de formas incompatibles.
    \item \textbf{Cálculo Numérico (Derivación e Integración) (3):} Valida la precisión de la derivación numérica frente a soluciones analíticas de referencia mediante el cálculo explícito de errores RMS y máximos con tolerancia. Adicionalmente, se verifica la consistencia de la integración numérica mediante la validación de la regla del trapecio.
    \item \textbf{Simulación y Reproducibilidad (1):} Verificación de la consistencia de los ejes temporales, de las gráficas y comprobación del comportamiento determinista del pipeline ante una semilla idéntica.
\end{itemize}
La ejecución completa de la suite de pruebas toma segundos en el entorno de desarrollo, garantizando un mecanismo de retroalimentación rápido y un comportamiento determinista en cualquier sistema externo.

\section{Reproducibilidad}
Para garantizar la ejecución en cualquier equipo externo sin transferir entornos virtuales ni modificar rutas, el flujo se ejecuta mediante los comandos públicos:
\begin{verbatim}
uv sync --locked
uv run ruff check .
uv run ruff format --check .
uv run pytest
uv run python scripts/reproduce.py
latexmk -pdf -interaction=nonstopmode -halt-on-error -outdir=build/reports reports/informe.tex
\end{verbatim}
Esto confirma que una segunda ejecución consecutiva reemplaza de forma controlada los archivos derivados sin duplicar registros ni alterar los resultados estocásticos.

\section{Limitaciones}
Más allá de las restricciones inherentes al uso de señales sintéticas —las cuales no contemplan perturbaciones severas o ruido realista complejo—, la principal limitación encontrada radica en la curva de aprendizaje y fricción operativa en el flujo de trabajo. 

El uso estricto de herramientas modernas de gestión de dependencias y empaquetado (como \texttt{uv}), la separación estricta de código bajo la carpeta \texttt{src/} y la configuración de pruebas automatizadas mediante \texttt{pytest} no siempre resultan fluidas al inicio. Esto genera barreras de adopción que exigen un rigor metodológico elevado y un dominio técnico previo que puede ralentizar el desarrollo inicial.

\section{Contribuciones}
Las responsabilidades del equipo se encuentran distribuidas y verificadas en el control de versiones: diseño de módulos de calidad, implementación de métodos numéricos, construcción de pruebas unitarias y automatización documental.

\end{document}